\documentclass[10pt,a4paper]{article}

\usepackage[margin=1.05in]{geometry}
\usepackage{setspace}
\usepackage{amsmath,amssymb,amsthm}
\usepackage{graphicx}
\usepackage{array}
\usepackage{caption}
\usepackage{enumitem}
\usepackage{titlesec}
\usepackage{microtype}
\usepackage{hyperref}
\usepackage{fancyhdr}

\hypersetup{colorlinks=true,linkcolor=black,urlcolor=black,citecolor=black}
\captionsetup{font=small,labelfont=bf,skip=4pt,justification=raggedright}

\titleformat{\section}{\normalsize\bfseries}{}{0em}{}
\titleformat{\subsection}{\small\bfseries}{}{0em}{}
\titleformat{\subsubsection}{\small\itshape}{}{0em}{}
\titlespacing*{\section}{0pt}{1.35em}{0.5em}
\titlespacing*{\subsection}{0pt}{0.95em}{0.35em}

\theoremstyle{plain}
\newtheorem{proposition}{Proposition}
\theoremstyle{definition}
\newtheorem{definition}{Definition}

\setstretch{1.14}
\setlist[itemize]{leftmargin=1.15em,itemsep=0.1em,topsep=0.2em}

\graphicspath{{drawings/}}

\pagestyle{fancy}
\fancyhf{}
\fancyhead[L]{\small\itshape Elements of Position}
\fancyhead[R]{\small\itshape V.\ Where 1 Is}
\fancyfoot[C]{\small\thepage}
\renewcommand{\headrulewidth}{0pt}

\begin{document}

\begin{center}
\vspace*{2.2em}
{\large Elements of Position}\\[1.6em]
{\LARGE\bfseries V.\ Where 1 Is}\\[2.2em]
{\normalsize Justin Erholtz}\\[0.35em]
{\small 12 September 2026}
\end{center}

\vspace{2.4em}

\begin{center}
{\large Where is \(1\)?}\\[0.35em]
{\large Point to it.}
\end{center}

\vspace{1.8em}

\noindent
This part does not replace I--IV.
It is a reflection, and it is also the last piece of the argument those parts were built to carry.

The first sentence of I was the question; II named $2$ from that $1$; III sat meetings on the same line; IV said the meetings that can be seated occupy that cut.
The question is still the question.

\vspace{1.4em}

\noindent
\textit{A note.}
Claims already proved in I--IV are used, not reopened.
The drawings are layers on the same architectural set, built by \texttt{Elements of Position - Drawings.py}.

\vspace{0.9em}
\noindent
\textit{A note on names used here.}
Names from I--IV stand.
HERE --- this pause, this finished count, this own-origin.
THERE --- the elected unit: the lock, the comparison that has arrived.

\newpage
\tableofcontents
\newpage

% =====================================================================
\section{The question, asked again}
% =====================================================================

Where is $1$?
Point to it.

You cannot, and not because the mark is missing.
One is not an object on the table.
It is a concept: a finished hold, the interior of which can be cut without end.

Between $0$ and $1$ the cut can be repeated without end.
That was already on the first page of I, and it was not decoration.
The endless is not waiting past the last integer.
The endless is already the interior of the unit.

\begin{figure}[ht]
\centering
\includegraphics[width=0.95\textwidth]{layer_d01_interior_endless.png}
\caption{Between $0$ and $1$ the marks do not finish.
The midpoint is the same cut asked again.
The endless is interior.}
\end{figure}

So $1$ is not a pebble.
It is the lock of a segment that already contains an infinity of cuts.
Normalize that interior to the unit --- that normalization \emph{is} $1$.

I called that object the closed unit; II asked it the other way and named $2$; III found meetings on its even cut; IV said there is no other primary line for a seated meeting to occupy.
None of that moved the question --- it only built the place the question was already pointing at.

% =====================================================================
\section{HERE and THERE}
% =====================================================================

$1$ is THERE.
We are HERE.

THERE is the elected unit: the comparison that has arrived, the lock $x=1/x$, the geometric mean of every inversion pair.
HERE is the pause on the street: this station, this finished integer, this own-origin.

Arrive at THERE and you are HERE --- same place, other writing.
Inversion says that in one stroke, but trace has to walk it.

\begin{figure}[ht]
\centering
\includegraphics[width=0.72\textwidth]{layer_d02_here_there.png}
\caption{$1$ is THERE.
The pause is HERE.
The walk that closes is the same point as the unit it elects.}
\end{figure}

A name without a walk is vacuum; a walk without an arrival is not a unit.
Privilege is treating the name as the seat.

IV already used those two respects as poles of one relation; this part only says what they do to the first sentence.
Get THERE and you are HERE.
The answer and the asker occupy the same mark, written twice.

% =====================================================================
\section{Every tick is the same tick}
% =====================================================================

Infinity is not a later integer.
It is not a mystical last stop.
It is not a destination the landscape can finish in advance so the object will already be complete.

Every tick is the same package.
Diameter $n$, own-origin $n/2$, inverse $1/n$, half-walk $n\pi/2$, even cut $1/2$.
Labels change.
The drawing does not.

\begin{figure}[ht]
\centering
\includegraphics[width=0.95\textwidth]{layer_d03_every_tick.png}
\caption{$n=2$, $n=8$, $n=16$.
Same package.
The midline does not move.}
\end{figure}

\begin{figure}[ht]
\centering
\includegraphics[width=0.95\textwidth]{layer_d05_zoom.png}
\caption{Zoom is the same cut asked again.
Scale is a ratio of parts.
The page can change size.
The ratio cannot.}
\end{figure}

Rate is that sentence with a tempo stuck on.
One day between counts, one year, one millisecond: still the same tick.
No tempo reaches a different kind of mark called $\infty$.
If you could arrive at infinity you would have arrived at another HERE, and you would not know the difference.
You do not know what you do not know until the next mark is on the table, and the next mark is always another finite tick.

Infinity is always THERE, never HERE: HERE is a pause, and THERE is the unit the pause elects.

That is why I refused a last index, why II refused a completed list of primes, why III refused a completed Fibonacci sequence, and why IV refused two centres at one pause.
The refusals are one refusal.
The tick does not change its kind.

% =====================================================================
\section{The other pole of $1$}
% =====================================================================

The map $x\mapsto 1/x$ sends $0$ toward $\infty$ and $\infty$ toward $0$, fixing $1$.
The partner of the pole $0$ is not a last index of the $n$-count.
It is the other writing of the lock.

\begin{figure}[ht]
\centering
\includegraphics[width=0.78\textwidth]{layer_d04_poles.png}
\caption{$0\leftrightarrow\infty$ through $1$.
Infinity is the other pole of the unit, not a later integer.}
\end{figure}

\begin{figure}[ht]
\centering
\includegraphics[width=0.95\textwidth]{layer_d06_inversion.png}
\caption{Inversion through the lock.
$2$ is already named from $1/2$.
The axis has no last index.}
\end{figure}

Usual attempts seat infinity as a completed HERE: all arrivals listed, all meetings listed, the product already infinite, the object already closed before the walk.
They put THERE in the place reserved for HERE.
Then they hunt leftovers off a cut of an object they declared finished before the count.

The order here is the other way: the endless is already in $(0,1)$, normalized to $1$.
It does not get to arrive as a last index --- it stays THERE.
Every actual tick is finite, and every finite tick has the same midline.

That is why infinity is not a special tick --- it is every tick.
It is HERE, written as THERE.

Vitruvius, opening the book on water, said the same thing in another room:
nothing suffers annihilation, but at dissolution there is a change, and things fall back to the essential element in which they were before.
Dissolution is not a vanishing --- it is a change of writing.
The thing falls back to what named it.

% =====================================================================
\section{The answer}
% =====================================================================

Where is $1$?

It is the closed unit.
It is the lock of inversion.
It is the geometric mean of every pair $\{x,1/x\}$.
It is the whole of which $1/2$ is the only primary part.
It is THERE.

We are HERE: this pause, this station, this finished count.
Get THERE and you are HERE.

Between $0$ and $1$ the marks do not finish, so the endless is not elsewhere.
Infinity is the interior of that unit, and the other pole of that unit, and every tick of that unit asked again.
It is not a destination.
It is not imported.
It does not hide a second cut.

II's other walk was that interior asked as a process, after arrival; III's meetings were marks on the even cut of that same unit; IV's midline was the statement that a seated meeting has nowhere else to sit.
The location did not move.

The first sentence asked a location.
The location is the arrival.
The arrival is the unit.
The unit was the question.

% =====================================================================
\section{What stands}
% =====================================================================

Forced, and already forced in I--IV:
$1$ as lock and as interior;
HERE as pause, THERE as elected unit;
every finished $n$ the same package;
zoom the same cut asked again;
$0\leftrightarrow\infty$ through $1$, not through a last index;
meetings that can be seated sit on $1/2$.

Named:
HERE, THERE.

Not granted:
infinity as a last tick;
THERE seated in the place reserved for HERE;
this reflection as optional colour.

The question is still the first sentence.
The answer is still the same unit.
Point to it.
You are pointing at the hold you already closed, and at the cut that hold still admits.
Both are $1$.

\end{document}
